\documentclass[11pt]{article}
\usepackage[margin=1in]{geometry}
\usepackage{amsmath,amssymb,amsthm,booktabs}
\usepackage[T1]{fontenc}
\usepackage{lmodern}
\usepackage[colorlinks=true,linkcolor=blue,citecolor=blue,urlcolor=blue]{hyperref}
\usepackage{microtype}
\newtheorem{theorem}{Theorem}
\newtheorem{lemma}{Lemma}
\newtheorem{proposition}{Proposition}
\newcommand{\sn}{\operatorname{sn}}
\newcommand{\cn}{\operatorname{cn}}
\newcommand{\dn}{\operatorname{dn}}
\newcommand{\Area}{\mathcal A}
\newcommand{\ii}{\mathrm i}
\title{A two-pole proof of the centered pedal-area invariant}
\author{Alec Kriebel\\\small Independent researcher\\
\small\href{https://orcid.org/0009-0001-9320-500X}{ORCID: 0009-0001-9320-500X}}
\date{16 October 2026}
\begin{document}
\maketitle
\begin{abstract}
We give a Jacobi elliptic-function proof that the signed area of a periodic
elliptic-billiard orbit times the signed area of its pedal polygon from the
common center is constant for least periods odd or divisible by four.
The same scoped theorem appears in Ferudun's contemporaneous preprint;
this is an attributed alternative proof, with no theorem, formula, method,
priority or independent-discovery claim. The proof identifies the pedal
area with a shifted cyclic $\dn$ sum by canceling apparent double poles
between incident area terms. A two-pole divisor then gives the invariant.
Primitive star orbits are included. We make the necessary least-period
convention explicit and give an exact repeated-orbit counterexample to an
unrestricted list-length formulation. Verification materials accompany this
AI-assisted, unrefereed research note.
\end{abstract}

\section{Statement, scope and attribution}
Fix $a>b>0$ and $0<\lambda<b^2$, and consider
\begin{equation}\label{ellipses}
 E:\frac{x^2}{a^2}+\frac{y^2}{b^2}=1,
 \qquad C:\frac{x^2}{a^2-\lambda}+\frac{y^2}{b^2-\lambda}=1.
\end{equation}
Thus $C$ is a proper, strictly nested confocal elliptical caustic. For an
ordered closed polygon $V=(V_j)_{j\in\mathbb Z/N\mathbb Z}$ set
\[
 \Area(V)=\frac12\sum_{j=0}^{N-1}\det(V_j,V_{j+1}).
\]
Let $P$ be a physical billiard orbit on $E$ tangent to $C$. Its ordered
center-pedal polygon consists of the perpendicular feet $Q_j$ from
$O=(0,0)$ to the \emph{complete original sidelines} $P_jP_{j+1}$.
Write $A=\Area(P)$ and $B=\Area(Q)$.

\begin{theorem}\label{main}
If the least physical period $N\geq3$ is odd or divisible by four, then
$AB$ is independent of phase in the fixed-caustic family
\eqref{ellipses}. Every admitted coprime winding number, including
primitive stars, and either orientation are allowed.
\end{theorem}

The observable, parity and center restriction are the experimentally
reported invariant $k_{203,b}$ of Reznik, Garcia and Koiller
\cite[Table~3]{RGK,RGKarxiv}. Their area convention is signed, and
their table does not literally specify least period. We state that
qualification explicitly: Section~\ref{repetition} shows that a version
based only on the length of a repeated vertex list is false. No
hyperbolic, collapsed or nonconfocal caustic, moving pedal point, clipped
segment projection, or unweighted filled-region area of a star is covered.

Ferudun \cite[Theorem~1]{Ferudun} proves exactly this qualified theorem,
including its star and repeated-list conventions. We credit that result
prominently and compare the trace formulas in Section~\ref{comparison}.
The present note records another proof organization and reproducibility
materials; it does not claim a new resolution of a presently open problem.
Stachel's canonical parametrization and the established meromorphic
methods of \cite{AST,Chavez} are prior mathematics.

\section{Canonical coordinates and the two areas}
Put
\[
 \alpha=\sqrt{a^2-\lambda},\quad
 \beta=\sqrt{b^2-\lambda},\quad
 k=\frac{\sqrt{\alpha^2-\beta^2}}{\alpha},\quad k'=\frac\beta\alpha.
\]
All Jacobi functions below have modulus $k\in(0,1)$. Let $K=K(k)$ and
$K'=K(k')$ be the complete elliptic integrals of the first kind.
Stachel \cite[Theorem~4.3 and (4.9)]{Stachel} gives
\begin{equation}\label{canonical}
 P(w)=(-a\sn w,b\cn w),\qquad
 a=\frac{\alpha\dn v}{\cn v},\quad b=\frac\beta{\cn v},\quad
 \delta=2v=\frac{4K\tau}{N},
\end{equation}
where $0<v<K$, $0<\tau<N/2$ and $\gcd(N,\tau)=1$, after choosing
the orientation. The vertices are $P(w+j\delta)$. Conversely these
parameters describe every proper pair \eqref{ellipses} admitting such a
periodic family: $\cn v=\beta/b$ has a unique solution in $(0,K)$,
and confocality gives the other relation in \eqref{canonical}.
The real map $P$ parametrizes $E$ once modulo $4K$, so the coprimality
condition is precisely the least-period condition. It does not restrict
the winding number to one.

The contact phase is useful to verify directly. The line through
$P(u-v)$ and $P(u+v)$ is
\begin{equation}\label{tangent}
 n(u)\cdot X=1,\qquad
 n(u)=\left(-\frac{\sn u}{\alpha},\frac{\cn u}{\beta}\right).
\end{equation}
It is tangent to $C$ at $(-\alpha\sn u,\beta\cn u)$. Substitution
of either endpoint using the addition formulas reduces the numerator to
$\sn^2u\dn^2v+\cn^2u=1-k^2\sn^2u\sn^2v$, the addition
denominator. The actual Euclidean center foot is therefore
\begin{equation}\label{Q}
 Q(u)=\frac{n(u)}{n(u)\cdot n(u)}
 =\left(-\frac{\alpha k'^2\sn u}{\dn^2u},
             \frac{\beta\cn u}{\dn^2u}\right).
\end{equation}
For real $u$, $n(u)\cdot n(u)=\dn^2u/\beta^2>0$. In complex
continuation the dot product is bilinear, with no conjugation.

Define
\begin{equation}\label{traces}
 S(w)=\sum_{j=0}^{N-1}\dn(w+j\delta),\qquad
 T(u)=\frac12\sum_{j=0}^{N-1}
           \det\bigl(Q(u+j\delta),Q(u+(j+1)\delta)\bigr).
\end{equation}
Then $B(w)=T(w+v)$. The addition identities give
\begin{align*}
 \det(P(u-v),P(u+v))
  &=\frac{2ab\sn v\cn v\dn u}{1-k^2\sn^2v\sn^2u},\\
 \dn(u-v)+\dn(u+v)
  &=\frac{2\dn v\dn u}{1-k^2\sn^2v\sn^2u}.
\end{align*}
After the shoelace factor $1/2$, summing counts each vertex $\dn$
term twice. Hence
\begin{equation}\label{A}
 A(w)=C_A S(w),\qquad C_A=\frac{ab\sn v\cn v}{\dn v}>0.
\end{equation}
These are signed determinant identities; no simplicity assumption on the
polygon has been made.

\section{The divisor of the cyclic delta-amplitude}
Set
\begin{equation}\label{reduction}
 m=\begin{cases}N&N\text{ odd},\\N/2&N\text{ even},\end{cases}
 \qquad \ell=\frac{2K}{m},\qquad p=\ii K'.
\end{equation}
The order of $\delta$ modulo $2K$ is $m$. Its $N$ residues have
multiplicity $N/m$, and the generating integer $\delta/\ell$ is
coprime to $m$. Reindexing and the $2K$ period of $\dn$ show that
$S$ is even and $\ell$-periodic. It has period $4p$ and antiperiod $2p$.
We use the standard poles, zeros and shifts in \cite[\S22.4]{DLMF};
in particular $\dn(z+2K)=\dn z$ and $\dn(z+2p)=-\dn z$.

\begin{lemma}\label{Slemma}
On $X=\mathbb C/(\ell\mathbb Z+4p\mathbb Z)$,
\[
 \operatorname{div}(S)
 =[p+\ell/2]+[3p+\ell/2]-[p]-[3p].
\]
Consequently $S(w)S(w+\ell/2)$ is a positive real constant for real $w$.
\end{lemma}
\begin{proof}
The poles of $\dn$ are simple and occur at
$2rK+(2s+1)p$. On $X$ the only possible poles of $S$ are $p$ and
$3p$. At $p$ exactly $N/m$ summands have a pole; their arguments
differ by real periods, so their residues agree and do not cancel.
Indeed $\dn(p+z)=-\ii/z+O(z)$, so the residue is $-\ii N/m$.
The antiperiod gives the opposite residue at $3p$.

At $z=p+\ell/2$, which is not a pole, the equality
$-z+\ell+2p=z$ combines evenness and the antiperiod to give
$S(z)=-S(z)$. The same holds at $3p+\ell/2$. These two distinct
zeros exhaust the zero divisor, since an elliptic function with two
simple poles has total zero order two. Both zeros are simple.
The divisor of $S(w)S(w+\ell/2)$ is therefore zero; compactness
makes the product constant. On the real axis every summand $\dn$ is
positive, so the constant is positive.
\end{proof}

\section{Cancellation in the pedal-area trace}
\begin{lemma}\label{Tlemma}
For a real constant $C_T$, one has $T(u)=C_T S(u+K)$. In fact
\begin{equation}\label{CT}
 C_T=\alpha^2\frac{\sn\delta(\dn\delta+k'^2)}{(1+\dn\delta)^2}>0.
\end{equation}
\end{lemma}
\begin{proof}
At a common pole of $\sn,\cn,\dn$, the numerators in \eqref{Q}
have order one and $\dn^2$ order two; $Q$ is holomorphic there and
vanishes. Its remaining possible poles are the simple zeros of $\dn$,
at $r=K+p$ and their translates. They have order at most two for $Q$.
The shift identities and parity give
\begin{equation}\label{even}
 Q(2r-z)=Q(z).
\end{equation}
Thus $Q(r+\varepsilon)$ has only even Laurent powers, with no
$\varepsilon^{-1}$ term.

At a phase $u+j\delta=r+\varepsilon$, the two incident summands
of the ordered area trace combine as
\begin{equation}\label{incident}
 \frac12\det\bigl(Q(r+\varepsilon),
            Q(r+\delta+\varepsilon)-Q(r-\delta+\varepsilon)\bigr).
\end{equation}
The neighbors are holomorphic because $0<\delta<2K$ is not a
real pole-to-pole shift. Their difference vanishes at $\varepsilon=0$
by \eqref{even}. Its order at least one reduces the apparent double
pole of \eqref{incident} to at most a simple pole. If $N$ is even,
the $N/m=2$ singular indices are separated by $m\geq2$; no edge has
two singular endpoints, so the incident grouping remains valid without
counting an edge twice. This includes $N=4$, when the regular
neighbors vanish. For $N=3$ there is only one singular index.

From \eqref{Q},
\[
 Q(u+2K)=-Q(u),\qquad
 Q(u+2p)=\operatorname{diag}(1,-1)Q(u).
\]
Determinants make $T$ periodic under $2K$ and antiperiodic under $2p$.
The closed cyclic sum is also $\delta$-periodic, hence $\ell$-periodic.
It is meromorphic on $X$, with only possible simple poles at $r$ and
$r+2p$. These are exactly the poles of $S(u+K)$, modulo the real
period $2K$. Match one residue by a scalar $C_T$; the antiperiod
matches the other. The difference is a holomorphic function on the
compact torus, hence constant, and its antiperiod forces it to be zero.

For completeness we evaluate that scalar. The quarter-period formulas
give
\[
 Q(r+x)=\frac\alpha k\frac{\cn x}{\sn^2x}(\dn x,\ii),
 \qquad U=\frac\alpha k(1,\ii),
 \qquad Q(r+\varepsilon)=U\varepsilon^{-2}+O(1).
\]
The neighbor difference in \eqref{incident} is
$2\varepsilon Q'(r+\delta)+O(\varepsilon^3)$; the residue of
$T$ is $(N/m)\det(U,Q'(r+\delta))$. The residue of $S(u+K)$
is $-\ii N/m$. With $q_z=Q_x+\ii Q_y$, the transformed formula
simplifies to
\[
 q_z(r+x)=-\frac{\alpha k\cn x}{1+\dn x},\qquad
 q_z'(r+x)=\frac{\alpha k\sn x(\dn x+k'^2)}{(1+\dn x)^2}.
\]
Therefore $C_T=\ii\det(U,Q'(r+\delta))=(\alpha/k)q_z'(r+\delta)$,
which is \eqref{CT}. Since $0<\delta<2K$, its real factors give
$C_T>0$. No division by a polygon area was used.
\end{proof}

\section{Parity and conclusion}
Equations \eqref{A} and Lemma~\ref{Tlemma} give
\begin{equation}\label{product}
 A(w)B(w)=C_A C_T S(w)S(w+K+v).
\end{equation}
If $N$ is odd, $(K+v)/\ell=N/2+\tau$ is a half-integer.
If $N$ is even it equals $(m+\tau)/2$. For $4\mid N$, $m$ is
even and $\tau$ odd, so this too is a half-integer. In both target
cases $K+v\equiv\ell/2\pmod\ell$, and Lemma~\ref{Slemma}
proves Theorem~\ref{main}.

For excluded primitive $N\equiv2\pmod4$, $m$ and $\tau$ are
odd, so the shift is an integer multiple of $\ell$. The product is
$C_A C_T S(w)^2$ and is genuinely nonconstant: $S$ has poles and
both coefficients are positive. Constancy on a real interval would
contradict the meromorphic identity theorem. Reversal negates $A$ and
$B$ and preserves their product. In the chosen orientation both areas
are positive, by \eqref{A} and \eqref{CT}. Taking absolute values of
these signed shoelace areas also preserves the product; this says
nothing about unweighted filled-in areas of stars.

For example the least-four caustic has
$\lambda=a^2b^2/(a^2+b^2)$. At its axial diamond,
$A=2ab$ and $B=4a^3b^3/(a^2+b^2)^2$, hence the invariant is
$8\lambda^2$. This value and low-period results are also in the
overlapping literature \cite{Ferudun,FerudunGeneral}.

\section{Why a repeated list does not change the period}\label{repetition}
An exact boundary check uses $a=2$, $b=1$ and $\lambda=4/9$.
The following two positive physical six-period orbits have that caustic:
\begin{align*}
 H={}&\bigl((2,0),(4/3,\sqrt5/3),(-4/3,\sqrt5/3),\\
     &\hspace{16mm}(-2,0),(-4/3,-\sqrt5/3),(4/3,-\sqrt5/3)\bigr),\\
 J={}&\bigl((0,1),(-4\sqrt2/3,1/3),(-4\sqrt2/3,-1/3),\\
     &\hspace{16mm}(0,-1),(4\sqrt2/3,-1/3),(4\sqrt2/3,1/3)\bigr).
\end{align*}
Direct substitution verifies ellipse membership, tangent discriminants
and forward specular reflection at every vertex; all six vertices are
distinct. Direct perpendicular projection and shoelace summation give
\[
\begin{array}{c|ccc}
 & A&B&AB\\\hline
 H&20\sqrt5/9&140\sqrt5/81&14000/729\\
 J&32\sqrt2/9&224\sqrt2/81&14336/729
\end{array}
\]
Double traversal multiplies each area by two and its product by four.
The two twelve-entry products differ by $448/243$, while their least
period stays six. Thus a list length divisible by four is insufficient.
Ferudun \cite[Section~6]{Ferudun} already gives the same obstruction
with another exact physical example; no new boundary phenomenon is
claimed here.

\section{Comparison and the documented chronology}\label{comparison}
Let $W(u)=a^2-(a^2-b^2)\sn^2u$. Ferudun's reflection-coordinate
telescope expresses $B$ as a constant times
$\sum_j\dn(w+j\delta)/W(w+j\delta)$ \cite[Lemma~3]{Ferudun}.
The ordinary Jacobi addition and quarter-period formulas give
\[
 \dn(u+K-v)+\dn(u+K+v)=\frac{2ab\dn u}{W(u)}.
\]
For instance, its addition denominator satisfies
\[
 \dn^2u-k^2\cn^2u\sn^2v=\cn^2v\,W(u)/\alpha^2,
\]
and
$k'\dn(v)\alpha^2/\cn^2v=ab$. Summing around the closed
orbit reindexes the two shifted terms, giving
\[
 \sum_j\frac{\dn(w+j\delta)}{W(w+j\delta)}
  =\frac1{ab}S(w+K+v).
\]
Thus our normalization, including the explicit coefficient \eqref{CT},
is already implied by that preprint and a standard addition identity.
Our direct-foot pole comparison is a different organization, not a
verified new formula or method. Ferudun's later arbitrary-point paper
also uses adjacent-foot pole cancellation \cite[Lemma~4.3]{FerudunGeneral}.
Earlier meromorphic cancellation and compactness arguments are in
\cite{AST,Chavez}.

The GitHub server record for our original submission, PR204, reports
creation on 1 October 2026 at 05:28:14 UTC, with the preserved proof
at commit \texttt{551a524e2db6} (full identifier archived) \cite{PR204}.
Ferudun's centered Zenodo record reports creation on the same date at
15:39:54.821511 UTC; DataCite records registration at 15:39:55 UTC.
These are contemporaneous dated records. They do not by themselves
establish discovery times, earliest historical visibility, priority,
dependence, independence or literal simultaneity. Our original proof
materials for the neighboring $k_{203,a}$ and $k_{203,b}$ efforts were
shared as author contributions, not independent reviews. The present
origin-only proof does not assume an arbitrary-point invariant.

\section{Verification and disclosure}
The accompanying archive preserves the original submission, source and
priority ledgers, independent low-period derivations, and computational
checks. Independent Euclidean shooting used 63 primitive shape/winding
cases and 378 phase/direction orbits; exact radical controls verify two
phases each for periods three and four and the six-period counterexample.
Those are finite evidence. High-precision computations are noninterval
diagnostics, not a proof of the universal theorem.

The historical author/reviewer scripts reproduce their reported outputs
with assertions enabled, but lose correctness enforcement under Python
optimization. Separate explicit-exception derivatives preserve their
valid outputs and reject genuine false mathematical identities in both
modes. Current numerical diagnostics also reject nonfinite values.
The archive distinguishes untouched historical evidence from corrected
live verification code. The analytic proof above, rather than counts or
floating-point agreement, establishes the universal result.

AI tools were used extensively in the original derivation, mathematical
auditing, calculations, software and manuscript preparation. Internal
adversarial AI reviews do not constitute conventional human peer review.
This is an unrefereed preprint and has no proof-assistant certification.
The author is responsible for its claims. No outside person was contacted
on behalf of this research.

\begin{thebibliography}{9}
\bibitem{RGK} D. Reznik, R. Garcia and J. Koiller,
\emph{Fifty new invariants of $N$-periodics in the elliptic billiard},
Arnold Math. J. \textbf{7} (2021), 341--355.
\href{https://doi.org/10.1007/s40598-021-00174-y}{doi:10.1007/s40598-021-00174-y}.
\bibitem{RGKarxiv} D. Reznik, R. Garcia and J. Koiller,
\emph{Eighty new invariants of $N$-periodics in the elliptic billiard},
\href{https://arxiv.org/abs/2004.12497v11}{arXiv:2004.12497v11} (2020).
\bibitem{Ferudun} A. Ferudun,
\emph{Center-Pedal Area Products in Elliptic Billiards},
version 1.0, 1 October 2026, unrefereed preprint.
\href{https://doi.org/10.5281/zenodo.23085904}{doi:10.5281/zenodo.23085904}.
\bibitem{Stachel} H. Stachel, \emph{On the motion of billiards in ellipses},
Eur. J. Math. \textbf{8} (2022), 1602--1622.
\href{https://doi.org/10.1007/s40879-021-00524-2}{doi:10.1007/s40879-021-00524-2}.
\bibitem{DLMF} NIST Digital Library of Mathematical Functions,
\href{https://dlmf.nist.gov/22.4}{\S22.4},
\href{https://dlmf.nist.gov/22.8}{\S22.8} and
\href{https://dlmf.nist.gov/22.13}{\S22.13}, version 1.2.8,
checked 16 October 2026.
\bibitem{AST} A. Akopyan, R. E. Schwartz and S. Tabachnikov,
\emph{Billiards in ellipses revisited},
\href{https://arxiv.org/abs/2001.02934v2}{arXiv:2001.02934v2} (2020).
\bibitem{Chavez} A. C. Ch\'avez-Caliz,
\emph{More about areas and centers of Poncelet polygons},
Arnold Math. J. \textbf{7} (2021), 91--106.
\href{https://doi.org/10.1007/s40598-020-00154-8}{doi:10.1007/s40598-020-00154-8}.
\bibitem{FerudunGeneral} A. Ferudun,
\emph{Pedal Areas of Periodic Elliptic Billiard Trajectories with Respect
to an Arbitrary Point}, version 1.0, 15 October 2026, unrefereed preprint.
\href{https://doi.org/10.5281/zenodo.23400541}{doi:10.5281/zenodo.23400541}.
\bibitem{PR204} A. Kriebel, \emph{5100012: prove the k203,b origin-pedal
area-product invariant}, original GitHub PR204, server creation
1 October 2026. \url{https://github.com/AlecKriebel/Math/pull/204}.
\end{thebibliography}
\end{document}
